\begin{proposition}[KT-ODD-FAMILY-RATIONAL-MIRROR]
The weight progression in KT-ODD-DISCRIMINANT-TWO-UNIT-FAMILY
already satisfies the original rational class-component vanishing
by the reflection argument of K17-REFLECTOR.  More precisely, if
$p\geq7$, $p\equiv3\pmod4$, and
$e=(p-3)/2+(p-1)t$ with $t\geq0$, then
\[
 A_p^{[1+e]}=A_p^{[-e]}=0,
 \qquad K_{2e}(\mathbf Z)\{p\}=0.                     \tag{ODR.1}
\]
This uses neither auxiliary field $E_s$ nor the third unit (OD.5).
\end{proposition}
\begin{proof}
Put $q=(p-1)/2$, $F=\mathbf Q(\mu_p)$, and
$K=\mathbf Q(\sqrt{-p})$.  The character $\omega^q$ is the
quadratic character of $K$, and $1+e\equiv q\pmod{p-1}$.
The imaginary quadratic class-number formula used in (OD.6), now
with discriminant $-p$, gives $0<h(K)\leq(p-1)/2<p$.
Thus the $p$-primary class group of $K$ is zero.

Let $H=\operatorname{Gal}(F/K)$, of order $q$ prime to $p$.
Extension of ideals and norm satisfy
$j\circ N=\sum_{\sigma\in H}\sigma$ on the class group of $F$.
For an $H$-invariant $p$-primary class $x$ this gives
$q x=jN(x)=0$, since $N(x)$ is a $p$-primary class of $K$.
Multiplication by $q$ is invertible on a finite $p$-group, so
$A_p^H=0$.  In particular $A_p^{[q]}=0$.

Apply the reflection half of the already proved K17-REFLECTOR
at $e_0=(p-3)/2$.  This is even and lies in $[2,p-3]$;
its mirror character is $\omega^{1+e_0}=\omega^q$.
It follows that $A_p^{[-e_0]}=0$.  The character is unchanged
on replacing $e_0$ by $e$, and the carried K17-K2-PERIODIC
identification gives the last assertion of (ODR.1).
\end{proof}

\paragraph{Scope of the new auxiliary construction.}
The odd-discriminant theorem supplies actual GD fields and two
specified unit conditions, with both local splitting choices.
It does not remove a previously uncontrolled rational component
in this particular progression: (ODR.1) is a direct consequence
of the already accepted quadratic class-number and reflection
inputs.  A proof of the third unit would still be a separate
arithmetic statement about the auxiliary fields.  No full
third-unit assertion, no uniform result for arbitrary weight
residues, and no new general rank-three exclusion follows here.
